% Options for packages loaded elsewhere
% Options for packages loaded elsewhere
\PassOptionsToPackage{unicode}{hyperref}
\PassOptionsToPackage{hyphens}{url}
\PassOptionsToPackage{dvipsnames,svgnames,x11names}{xcolor}
%
\documentclass[
  spanish,
  letterpaper,
  DIV=11,
  numbers=noendperiod]{scrartcl}
\usepackage{xcolor}
\usepackage{amsmath,amssymb}
\setcounter{secnumdepth}{5}
\usepackage{iftex}
\ifPDFTeX
  \usepackage[T1]{fontenc}
  \usepackage[utf8]{inputenc}
  \usepackage{textcomp} % provide euro and other symbols
\else % if luatex or xetex
  \usepackage{unicode-math} % this also loads fontspec
  \defaultfontfeatures{Scale=MatchLowercase}
  \defaultfontfeatures[\rmfamily]{Ligatures=TeX,Scale=1}
\fi
\usepackage{lmodern}
\ifPDFTeX\else
  % xetex/luatex font selection
\fi
% Use upquote if available, for straight quotes in verbatim environments
\IfFileExists{upquote.sty}{\usepackage{upquote}}{}
\IfFileExists{microtype.sty}{% use microtype if available
  \usepackage[]{microtype}
  \UseMicrotypeSet[protrusion]{basicmath} % disable protrusion for tt fonts
}{}
\makeatletter
\@ifundefined{KOMAClassName}{% if non-KOMA class
  \IfFileExists{parskip.sty}{%
    \usepackage{parskip}
  }{% else
    \setlength{\parindent}{0pt}
    \setlength{\parskip}{6pt plus 2pt minus 1pt}}
}{% if KOMA class
  \KOMAoptions{parskip=half}}
\makeatother
% Make \paragraph and \subparagraph free-standing
\makeatletter
\ifx\paragraph\undefined\else
  \let\oldparagraph\paragraph
  \renewcommand{\paragraph}{
    \@ifstar
      \xxxParagraphStar
      \xxxParagraphNoStar
  }
  \newcommand{\xxxParagraphStar}[1]{\oldparagraph*{#1}\mbox{}}
  \newcommand{\xxxParagraphNoStar}[1]{\oldparagraph{#1}\mbox{}}
\fi
\ifx\subparagraph\undefined\else
  \let\oldsubparagraph\subparagraph
  \renewcommand{\subparagraph}{
    \@ifstar
      \xxxSubParagraphStar
      \xxxSubParagraphNoStar
  }
  \newcommand{\xxxSubParagraphStar}[1]{\oldsubparagraph*{#1}\mbox{}}
  \newcommand{\xxxSubParagraphNoStar}[1]{\oldsubparagraph{#1}\mbox{}}
\fi
\makeatother


\usepackage{longtable,booktabs,array}
\newcounter{none} % for unnumbered tables
\usepackage{calc} % for calculating minipage widths
% Correct order of tables after \paragraph or \subparagraph
\usepackage{etoolbox}
\makeatletter
\patchcmd\longtable{\par}{\if@noskipsec\mbox{}\fi\par}{}{}
\makeatother
% Allow footnotes in longtable head/foot
\IfFileExists{footnotehyper.sty}{\usepackage{footnotehyper}}{\usepackage{footnote}}
\makesavenoteenv{longtable}
\usepackage{graphicx}
\makeatletter
\newsavebox\pandoc@box
\newcommand*\pandocbounded[1]{% scales image to fit in text height/width
  \sbox\pandoc@box{#1}%
  \Gscale@div\@tempa{\textheight}{\dimexpr\ht\pandoc@box+\dp\pandoc@box\relax}%
  \Gscale@div\@tempb{\linewidth}{\wd\pandoc@box}%
  \ifdim\@tempb\p@<\@tempa\p@\let\@tempa\@tempb\fi% select the smaller of both
  \ifdim\@tempa\p@<\p@\scalebox{\@tempa}{\usebox\pandoc@box}%
  \else\usebox{\pandoc@box}%
  \fi%
}
% Set default figure placement to htbp
\def\fps@figure{htbp}
\makeatother


% definitions for citeproc citations
\NewDocumentCommand\citeproctext{}{}
\NewDocumentCommand\citeproc{mm}{%
  \begingroup\def\citeproctext{#2}\cite{#1}\endgroup}
\makeatletter
 % allow citations to break across lines
 \let\@cite@ofmt\@firstofone
 % avoid brackets around text for \cite:
 \def\@biblabel#1{}
 \def\@cite#1#2{{#1\if@tempswa , #2\fi}}
\makeatother
\newlength{\cslhangindent}
\setlength{\cslhangindent}{1.5em}
\newlength{\csllabelwidth}
\setlength{\csllabelwidth}{3em}
\newenvironment{CSLReferences}[2] % #1 hanging-indent, #2 entry-spacing
 {\begin{list}{}{%
  \setlength{\itemindent}{0pt}
  \setlength{\leftmargin}{0pt}
  \setlength{\parsep}{0pt}
  % turn on hanging indent if param 1 is 1
  \ifodd #1
   \setlength{\leftmargin}{\cslhangindent}
   \setlength{\itemindent}{-1\cslhangindent}
  \fi
  % set entry spacing
  \setlength{\itemsep}{#2\baselineskip}}}
 {\end{list}}
\usepackage{calc}
\newcommand{\CSLBlock}[1]{\hfill\break\parbox[t]{\linewidth}{\strut\ignorespaces#1\strut}}
\newcommand{\CSLLeftMargin}[1]{\parbox[t]{\csllabelwidth}{\strut#1\strut}}
\newcommand{\CSLRightInline}[1]{\parbox[t]{\linewidth - \csllabelwidth}{\strut#1\strut}}
\newcommand{\CSLIndent}[1]{\hspace{\cslhangindent}#1}

\ifLuaTeX
\usepackage[bidi=basic,shorthands=off]{babel}
\else
\usepackage[bidi=default,shorthands=off]{babel}
\fi
\ifLuaTeX
  \usepackage{selnolig} % disable illegal ligatures
\fi


\setlength{\emergencystretch}{3em} % prevent overfull lines

\providecommand{\tightlist}{%
  \setlength{\itemsep}{0pt}\setlength{\parskip}{0pt}}



 


\KOMAoption{captions}{tableheading}
\makeatletter
\@ifpackageloaded{caption}{}{\usepackage{caption}}
\AtBeginDocument{%
\ifdefined\contentsname
  \renewcommand*\contentsname{Tabla de contenidos}
\else
  \newcommand\contentsname{Tabla de contenidos}
\fi
\ifdefined\listfigurename
  \renewcommand*\listfigurename{Listado de Figuras}
\else
  \newcommand\listfigurename{Listado de Figuras}
\fi
\ifdefined\listtablename
  \renewcommand*\listtablename{Listado de Tablas}
\else
  \newcommand\listtablename{Listado de Tablas}
\fi
\ifdefined\figurename
  \renewcommand*\figurename{Figura}
\else
  \newcommand\figurename{Figura}
\fi
\ifdefined\tablename
  \renewcommand*\tablename{Tabla}
\else
  \newcommand\tablename{Tabla}
\fi
}
\@ifpackageloaded{float}{}{\usepackage{float}}
\floatstyle{ruled}
\@ifundefined{c@chapter}{\newfloat{codelisting}{h}{lop}}{\newfloat{codelisting}{h}{lop}[chapter]}
\floatname{codelisting}{Listado}
\newcommand*\listoflistings{\listof{codelisting}{Listado de Listados}}
\makeatother
\makeatletter
\makeatother
\makeatletter
\@ifpackageloaded{caption}{}{\usepackage{caption}}
\@ifpackageloaded{subcaption}{}{\usepackage{subcaption}}
\makeatother
\usepackage{bookmark}
\IfFileExists{xurl.sty}{\usepackage{xurl}}{} % add URL line breaks if available
\urlstyle{same}
\hypersetup{
  pdftitle={Brecha de detección de demencia y priorización de establecimientos en la red pública chilena: un análisis multinivel de los registros REM 2025},
  pdfauthor={Javier Vera Bravo},
  pdflang={es},
  colorlinks=true,
  linkcolor={blue},
  filecolor={Maroon},
  citecolor={Blue},
  urlcolor={Blue},
  pdfcreator={LaTeX via pandoc}}


\title{Brecha de detección de demencia y priorización de
establecimientos en la red pública chilena: un análisis multinivel de
los registros REM 2025}
\author{Javier Vera Bravo}
\date{2026-06-25}
\begin{document}
\maketitle

\renewcommand*\contentsname{Tabla de contenidos}
{
\setcounter{tocdepth}{3}
\tableofcontents
}

\subsection*{Resumen}\label{resumen}
\addcontentsline{toc}{subsection}{Resumen}

\textbf{Antecedentes.} La demencia tiene una frecuencia conocida por
edad, pero la fracción efectivamente detectada y bajo seguimiento en los
sistemas públicos es baja y heterogénea. En Chile no existían
estimaciones comunales de la brecha de detección basadas en el registro
administrativo con un denominador de prevalencia acompañado de
incertidumbre.

\textbf{Objetivos.} (1) Cuantificar la brecha de detección de demencia
en personas de 65 años o más en la red pública chilena y su gradiente
territorial; (2) determinar en qué nivel del sistema (establecimiento,
comuna, región) reside la variación del registro; y (3) derivar, a
partir de ese análisis, un indicador de priorización de establecimientos
para focalizar la intervención.

\textbf{Métodos.} Estudio ecológico de fuentes administrativas (Series
REM 2025, DEIS-MINSAL). El numerador es la población bajo control de
salud mental por demencia (65+, corte de diciembre); el denominador, la
prevalencia esperada por banda etaria aplicada a las proyecciones de
población del INE y escalada a la población inscrita en FONASA. La
incertidumbre del denominador se propagó por simulación. La partición de
la varianza del registro se estimó con un modelo logístico multinivel de
interceptos aleatorios (\texttt{glmmTMB}), reportando VPC y \emph{Median
Odds Ratio} (MOR). El puntaje de priorización combina el efecto
aleatorio por establecimiento con el volumen de población inscrita.

\textbf{Resultados.} La brecha nacional fue de 93.7 \% (sensibilidad 92
\%--94 \% según la curva de prevalencia). La detección urbana casi
duplica la rural. La variación del registro es predominantemente
institucional: el establecimiento explica el 53.2\% de la varianza (MOR
15.96).

\textbf{Conclusiones.} La mayor parte de la demencia esperada no está
bajo seguimiento de la salud mental pública, con una desventaja rural
marcada; la palanca de intervención es el establecimiento, coordinado
por Servicio de Salud. El indicador de priorización traduce el hallazgo
en una lista accionable de centros.

\emph{Palabras clave:} demencia; brecha de detección; epidemiología de
servicios; análisis multinivel; atención primaria; Chile.

\section{Introducción}\label{introducciuxf3n}

La demencia es una de las principales causas de dependencia en personas
mayores y su prevalencia crece de forma aproximadamente exponencial con
la edad (Fuentes y Albala 2014; {Llibre Rodriguez et~al.} 2008). A
escala global, una fracción mayoritaria de los casos no está
diagnosticada: el World Alzheimer Report 2021 estima cerca de tres de
cada cuatro casos sin diagnóstico, proporción que aumenta en países de
ingreso bajo y medio (Alzheimer's Disease International 2021). La
distancia entre la prevalencia esperada y los casos efectivamente
identificados ---la \emph{brecha de detección}--- es, por tanto, un
problema sistémico y no anecdótico (GBD 2019 Dementia Forecasting
Collaborators 2022).

En Chile, la prevalencia de demencia en población de 60 años o más se ha
estimado en torno al 7\% a partir de una encuesta nacional por tamizaje
(Fuentes y Albala 2014), y el Ministerio de Salud ha incorporado un
indicador oficial de cobertura que compara las personas bajo control por
demencia contra las esperadas por prevalencia (Ministerio de Salud de
Chile 2024). Sin embargo, hasta donde alcanza esta revisión, no se había
cuantificado esa brecha a nivel \textbf{comunal} usando el registro
administrativo REM contra un denominador de prevalencia con
incertidumbre propagada, ni se había determinado en qué nivel del
sistema reside la variación del registro.

Este trabajo aborda ese vacío con tres objetivos: cuantificar la brecha
y su gradiente territorial; particionar la varianza del registro entre
establecimiento, comuna y región; y derivar un indicador de priorización
de establecimientos que traduzca el hallazgo en acción de gestión.

\section{Métodos}\label{muxe9todos}

\subsection{Diseño y fuentes de
datos}\label{diseuxf1o-y-fuentes-de-datos}

Estudio ecológico de corte transversal sobre registros administrativos.
La fuente primaria son las Series REM 2025 del Departamento de
Estadísticas e Información de Salud (DEIS-MINSAL), datos preliminares.
Se complementan con las proyecciones de población del INE (base 2017),
los inscritos validados de FONASA en atención primaria, y el maestro de
establecimientos de salud (datos.gob.cl, licencia CC0). El cruce
territorial REM--INE y REM--maestro alcanzó el 100\% de las comunas y
establecimientos.

\subsection{Definición de caso}\label{definiciuxf3n-de-caso}

El universo de demencia se definió mediante un \emph{crosswalk} curado y
verificado código a código contra los diccionarios y manuales REM. El
numerador de la brecha es la población bajo control de salud mental por
``Demencias (incluye Alzheimer)'' en personas de 65 años o más, sumando
atención primaria (\texttt{P6222300}) y especialidad
(\texttt{P6223310}), verificadas como niveles de atención distintos del
mismo diagnóstico. Se utilizó un solo corte (diciembre), por tratarse de
un \emph{stock} (prevalencia), evitando el doble conteo de los dos
cortes semestrales.

\subsection{Denominador y triangulación de la
prevalencia}\label{denominador-y-triangulaciuxf3n-de-la-prevalencia}

La demencia esperada por comuna se construyó como

\[E_c = \kappa_c \sum_b P_{cb}\,\pi_b,\]

donde \(P_{cb}\) es la población de la comuna \(c\) en la banda etaria
\(b\) (INE), \(\pi_b\) la prevalencia por banda (Fuentes y Albala 2014)
y \(\kappa_c\) la cobertura pública (inscritos FONASA 60+ sobre
población INE 60+). El uso de la población pública corrige el sesgo de
contabilizar a beneficiarios de seguros privados que no usan la red
pública.

Dado que la prevalencia proviene de una sola encuesta por tamizaje, su
efecto sobre la brecha se evaluó por triangulación: la curva original,
una versión suavizada con monotonía impuesta, y un espacio para una
curva externa modelada (GBD 2019 Dementia Forecasting Collaborators
2022). La incertidumbre se propagó por simulación de Monte Carlo sobre
la prevalencia por banda, reportando un intervalo del 95\% para la
brecha nacional.

\subsection{Análisis espacial}\label{anuxe1lisis-espacial}

La autocorrelación espacial de la tasa de detección comunal se evaluó
con el índice de Moran global y estadísticos locales (LISA), con
contigüidad tipo \emph{queen} y pesos estandarizados por fila. El
análisis es ecológico y está sujeto al problema de la unidad areal
modificable (MAUP).

\subsection{Modelo multinivel}\label{modelo-multinivel}

La variación del registro se modeló como la probabilidad de que un
establecimiento de atención primaria activo registre actividad de
demencia en un mes dado (desenlace binario; panel establecimiento ×
mes). Se ajustó un modelo logístico de interceptos aleatorios de tres
niveles (Snijders y Bosker 2012; {Brooks et~al.} 2017):

\[\operatorname{logit}\big(\Pr(Y_{ijk}=1)\big) = \beta_0 + u_k + u_{jk} + u_{ijk}, \qquad u \sim \mathcal{N}(0,\sigma^2),\]

con \(k\) región, \(jk\) comuna e \(ijk\) establecimiento. La proporción
de varianza por nivel (VPC) se calculó en la escala latente, fijando la
varianza de nivel inferior en \(\pi^2/3\) (Goldstein et~al. 2002):

\[\text{VPC}_{\text{nivel}} = \frac{\sigma^2_{\text{nivel}}}{\sum_l \sigma^2_l + \pi^2/3}.\]

Como medida robusta a esa convención se reportó el \emph{Median Odds
Ratio} ({Merlo et~al.} 2006):

\[\text{MOR} = \exp\!\big(\sqrt{2\,\sigma^2}\,\Phi^{-1}(0{,}75)\big), \qquad \Phi^{-1}(0{,}75)\approx 0{,}6745.\]

Se ajustó un modelo alternativo reemplazando región por Servicio de
Salud. Ambos modelos son especificaciones no anidadas; el VPC, al ser un
modelo no condicionado, indica \emph{dónde} vive la variación, no
\emph{por qué} (Mood 2010). La convergencia se verificó en cada ajuste.

\subsection{Puntaje de priorización de
establecimientos}\label{puntaje-de-priorizaciuxf3n-de-establecimientos}

El indicador operativo de la herramienta de gestión combina, por
establecimiento \(i\):

\[S_i = 100 \cdot \tilde d_i \cdot \tilde v_i \cdot w_i, \qquad d_i = \max(0,\,-\hat u_i), \quad v_i = \log(1+\text{inscritos}_i),\]

donde \(\hat u_i\) es el efecto aleatorio del establecimiento (déficit
de registro ajustado por su contexto comunal y regional), \(v_i\) el
volumen de población expuesta, \(\tilde{\cdot}\) la normalización
min-max, y \(w_i\) un ponderador de equidad territorial (realce rural).
El puntaje prioriza centros que podrían registrar mucho más y atienden a
mucha gente. El desenlace es \textbf{visibilidad administrativa del
registro, no calidad clínica}.

\subsection{Reproducibilidad y gobernanza de
datos}\label{reproducibilidad-y-gobernanza-de-datos}

El análisis es reproducible de punta a punta
(\texttt{source("R/10\_run\_all.R")}), con entorno computacional fijado
(\texttt{renv}). Los datos crudos no se redistribuyen y conservan su
licencia de origen. Se aplicó un resguardo de celdas pequeñas (supresión
de tasas con menos de cinco personas) para evitar reidentificación; las
desagregaciones de pueblos originarios y migrantes se reportan solo
agregadas. El estudio usa datos administrativos agregados y públicos.

\section{Resultados}\label{resultados}

\subsection{Brecha de detección y sensibilidad del
denominador}\label{brecha-de-detecciuxf3n-y-sensibilidad-del-denominador}

La brecha nacional de detección fue de 93.7 \% (tasa de detección 6.3
\%). La Tabla~\ref{tbl-sens} muestra su robustez frente a la curva de
prevalencia: el cambio entre la curva cruda y la suavizada es de cerca
de un punto porcentual y el intervalo de incertidumbre se mantiene
estrecho.

{

\begin{longtable}[]{@{}lrll@{}}

\caption{\label{tbl-sens}Brecha de detección país bajo distintas curvas
de prevalencia, con intervalo del 95\%.}

\tabularnewline

\toprule\noalign{}
Fuente & Esperados 65+ & Brecha & IC 95\% \\
\midrule\noalign{}
\endhead
\bottomrule\noalign{}
\endlastfoot
10\_66\_raw & 206725 & 93.7 \% & 92.7 \%--94.5 \% \\
10\_66\_suave & 178712 & 92.7 \% & 91.6 \%--93.6 \% \\

\end{longtable}

}

\subsection{Gradiente territorial}\label{gradiente-territorial}

La detección urbana casi duplica la rural (Tabla~\ref{tbl-zona}). Lo
rural acumula una doble desventaja: mayor prevalencia esperada y mayor
subregistro. El análisis espacial mostró autocorrelación significativa
con el denominador público (Moran's I = 0,23), con un conglomerado de
mejor detección en el sur austral, encabezado por comunas sede de
dispositivos del Plan Nacional de Demencia.

{

\begin{longtable}[]{@{}lll@{}}

\caption{\label{tbl-zona}Brecha de detección por tipo de comuna.}

\tabularnewline

\toprule\noalign{}
Zona & Tasa de detección & Brecha \\
\midrule\noalign{}
\endhead
\bottomrule\noalign{}
\endlastfoot
mixto & 4.4 \% & 95.6 \% \\
rural & 4.9 \% & 95.1 \% \\
urbano & 8.5 \% & 91.5 \% \\

\end{longtable}

}

\subsection{¿Dónde vive la variación del
registro?}\label{duxf3nde-vive-la-variaciuxf3n-del-registro}

La partición de la varianza es nítida (Tabla~\ref{tbl-ml}): el
establecimiento concentra la mayor parte de la variación, con un MOR que
indica que dos centros cualesquiera de la misma comuna difieren
marcadamente en sus probabilidades de registrar. En el modelo
alternativo, el Servicio de Salud pesa más que la región geográfica.

{

\begin{longtable}[]{@{}lrrr@{}}

\caption{\label{tbl-ml}Partición de varianza del registro (modelo
región/comuna/establecimiento).}

\tabularnewline

\toprule\noalign{}
Nivel & Varianza & VPC (\%) & MOR \\
\midrule\noalign{}
\endhead
\bottomrule\noalign{}
\endlastfoot
region.(Intercept) & 1.331 & 8.4 & 3.01 \\
comuna.(Intercept) & 2.795 & 17.6 & 4.93 \\
id.(Intercept) & 8.435 & 53.2 & 15.96 \\

\end{longtable}

}

\subsection{Cascada de atención}\label{cascada-de-atenciuxf3n}

Los registros por etapa no forman un embudo (Tabla~\ref{tbl-casc}): el
registro de sospecha y diagnóstico es minúsculo frente al ingreso y el
bajo control, lo que indica que el primer escalón de la ruta casi no se
anota. Es un resultado limitante: el REM, tal como está, no permite
reconstruir la trayectoria individual.

{

\begin{longtable}[]{@{}
  >{\raggedright\arraybackslash}p{(\linewidth - 8\tabcolsep) * \real{0.6875}}
  >{\raggedleft\arraybackslash}p{(\linewidth - 8\tabcolsep) * \real{0.0750}}
  >{\raggedleft\arraybackslash}p{(\linewidth - 8\tabcolsep) * \real{0.0875}}
  >{\raggedleft\arraybackslash}p{(\linewidth - 8\tabcolsep) * \real{0.0750}}
  >{\raggedleft\arraybackslash}p{(\linewidth - 8\tabcolsep) * \real{0.0750}}@{}}

\caption{\label{tbl-casc}Registros por etapa de la cascada (no anidan:
lógicas y unidades distintas).}

\tabularnewline

\toprule\noalign{}
\begin{minipage}[b]{\linewidth}\raggedright
Etapa
\end{minipage} & \begin{minipage}[b]{\linewidth}\raggedleft
País
\end{minipage} & \begin{minipage}[b]{\linewidth}\raggedleft
Urbano
\end{minipage} & \begin{minipage}[b]{\linewidth}\raggedleft
Rural
\end{minipage} & \begin{minipage}[b]{\linewidth}\raggedleft
Mixto
\end{minipage} \\
\midrule\noalign{}
\endhead
\bottomrule\noalign{}
\endlastfoot
1. Sospecha (A06 consultorías) & 325 & 57 & 89 & 49 \\
2. Diagnóstico (A06 consultorías) & 467 & 98 & 114 & 72 \\
3. Ingreso SM (A05, flujo anual) & 15117 & 3645 & 1243 & 2390 \\
4. Bajo control (P6, stock dic) & 14514 & 3776 & 1710 & 2065 \\
5. Domiciliaria dep. severa c/demencia (P3, stock dic) & 23209 & 5177 &
3009 & 4461 \\

\end{longtable}

}

\subsection{Priorización de
establecimientos}\label{priorizaciuxf3n-de-establecimientos}

A partir del efecto aleatorio por establecimiento se priorizaron los
2.027 centros de atención primaria activos. La Tabla~\ref{tbl-prio}
resume la distribución por nivel de prioridad; la herramienta
interactiva permite a cada Servicio de Salud explorar su red y exportar
la lista accionable.

{

\begin{longtable}[]{@{}lr@{}}

\caption{\label{tbl-prio}Distribución de establecimientos por nivel de
prioridad de intervención.}

\tabularnewline

\toprule\noalign{}
Prioridad & Establecimientos \\
\midrule\noalign{}
\endhead
\bottomrule\noalign{}
\endlastfoot
Alta & 338 \\
Media & 333 \\
Baja & 1356 \\

\end{longtable}

}

\section{Discusión}\label{discusiuxf3n}

La magnitud de la brecha es consistente con la literatura internacional
sobre subdiagnóstico de demencia (Alzheimer's Disease International
2021) y con el orden de magnitud del problema en países de ingreso medio
({Llibre Rodriguez et~al.} 2008). El aporte de este trabajo no es la
cifra global sino su \textbf{desagregación territorial} y, sobre todo,
la \textbf{localización de la variación}: al mostrar que más de la mitad
de la varianza del registro reside en el establecimiento ({Merlo et~al.}
2006), el análisis reorienta la intervención desde la campaña regional
genérica hacia la gestión centro a centro, coordinada por Servicio de
Salud ---la unidad administrativa que el modelo alternativo identifica
como más relevante que la región geográfica.

El puntaje de priorización operacionaliza ese hallazgo. Al basarse en el
efecto aleatorio del establecimiento, aísla el componente institucional
---el intervenible--- del contexto comunal y regional, y lo pondera por
el volumen de población expuesta. Coincide además con el marco de
monitoreo oficial (Ministerio de Salud de Chile 2024), extendiéndolo a
una resolución accionable por la gestión.

La triangulación del denominador es relevante metodológicamente: muestra
que, aunque la prevalencia proviene de una fuente única y antigua, la
conclusión sobre la magnitud de la brecha es robusta a esa elección. La
distorsión introducida por contabilizar a la población de seguros
privados, en cambio, sí afectaba materialmente tanto el nivel de la
brecha como la señal territorial, que emergió solo al usar el
denominador público.

\section{Limitaciones}\label{limitaciones}

El numerador mide bajo control de salud mental, no toda la atención de
demencia: la brecha cuantifica la visibilidad de ese subsistema, no la
demencia sin ningún contacto con el sistema. El subregistro rural infla
la brecha rural, de modo que parte del gradiente es falta de registro
más que falta de detección. El denominador descansa en una encuesta de
prevalencia por tamizaje de 2010, sensible a la escolaridad, lo que
confunde parcialmente el exceso rural con educación; la triangulación
atenúa pero no elimina esta limitación. El análisis es ecológico y está
sujeto al MAUP. El desenlace del modelo multinivel es visibilidad
administrativa, no calidad clínica. Las desagregaciones de pueblos
originarios y migrantes son marginales independientes y se interpretan
con cautela.

\section{Conclusión}\label{conclusiuxf3n}

La mayor parte de la demencia esperada en personas mayores no está bajo
seguimiento de la salud mental pública chilena, con una desventaja rural
marcada y una variación que reside predominantemente en el
establecimiento. La herramienta de priorización derivada traduce este
diagnóstico en una lista de centros a intervenir, recalculable
periódicamente con el mismo código, como insumo para la gestión de la
red.

\section*{Declaraciones}\label{declaraciones}
\addcontentsline{toc}{section}{Declaraciones}

\textbf{Uso de inteligencia artificial.} El proyecto se desarrolló con
asistencia de IA para programación, análisis y redacción, bajo revisión
y decisión humanas; la responsabilidad por el contenido y las
conclusiones es del autor. \textbf{Conflictos de interés.} El autor
declara no tener conflictos de interés. \textbf{Disponibilidad.} Código
y productos en el repositorio del proyecto; datos crudos disponibles en
sus fuentes originales (DEIS, INE, FONASA, datos.gob.cl).

\section*{Referencias}\label{referencias}
\addcontentsline{toc}{section}{Referencias}

\protect\phantomsection\label{refs}
\begin{CSLReferences}{1}{1}
\bibitem[\citeproctext]{ref-adi2021}
Alzheimer's Disease International. 2021. \emph{World Alzheimer Report
2021: Journey through the diagnosis of dementia}. Alzheimer's Disease
International.
\url{https://www.alzint.org/resource/world-alzheimer-report-2021/}.

\bibitem[\citeproctext]{ref-brooks2017}
{Brooks, Mollie E., Kasper Kristensen, Koen J. van Benthem, et~al.}
2017. {«glmmTMB balances speed and flexibility among packages for
zero-inflated generalized linear mixed models»}. \emph{The R Journal} 9
(2): 378-400. \url{https://doi.org/10.32614/RJ-2017-066}.

\bibitem[\citeproctext]{ref-fuentes2014}
Fuentes, Patricio, y Cecilia Albala. 2014. {«An update on aging and
dementia in Chile»}. \emph{Dementia \& Neuropsychologia} 8 (4): 317-22.
\url{https://doi.org/10.1590/S1980-57642014DN84000003}.

\bibitem[\citeproctext]{ref-gbd2022}
GBD 2019 Dementia Forecasting Collaborators. 2022. {«Estimation of the
global prevalence of dementia in 2019 and forecasted prevalence in 2050:
an analysis for the Global Burden of Disease Study 2019»}. \emph{The
Lancet Public Health} 7 (2): e105-25.
\url{https://doi.org/10.1016/S2468-2667(21)00249-8}.

\bibitem[\citeproctext]{ref-goldstein2002}
Goldstein, Harvey, William Browne, y Jon Rasbash. 2002. {«Partitioning
variation in multilevel models»}. \emph{Understanding Statistics} 1 (4):
223-31. \url{https://doi.org/10.1207/S15328031US0104_02}.

\bibitem[\citeproctext]{ref-llibre2008}
{Llibre Rodriguez, Juan J., Cleusa P. Ferri, Daisy Acosta, et~al.} 2008.
{«Prevalence of dementia in Latin America, India, and China: a
population-based cross-sectional survey»}. \emph{The Lancet} 372 (9637):
464-74. \url{https://doi.org/10.1016/S0140-6736(08)61002-8}.

\bibitem[\citeproctext]{ref-merlo2006}
{Merlo, Juan, Basile Chaix, Henrik Ohlsson, et~al.} 2006. {«A brief
conceptual tutorial of multilevel analysis in social epidemiology: using
measures of clustering in multilevel logistic regression to investigate
contextual phenomena»}. \emph{Journal of Epidemiology \& Community
Health} 60 (4): 290-97. \url{https://doi.org/10.1136/jech.2004.029454}.

\bibitem[\citeproctext]{ref-minsal2024}
Ministerio de Salud de Chile. 2024. \emph{Mesa Asesora de Demencia:
recomendaciones a la autoridad}. Ministerio de Salud, DIPRECE.
\url{https://diprece.minsal.cl/}.

\bibitem[\citeproctext]{ref-mood2010}
Mood, Carina. 2010. {«Logistic regression: Why we cannot do what we
think we can do, and what we can do about it»}. \emph{European
Sociological Review} 26 (1): 67-82.
\url{https://doi.org/10.1093/esr/jcp006}.

\bibitem[\citeproctext]{ref-snijders2012}
Snijders, Tom A. B., y Roel J. Bosker. 2012. \emph{Multilevel Analysis:
An Introduction to Basic and Advanced Multilevel Modeling}. 2nd ed.
Sage.

\end{CSLReferences}




\end{document}
